% Proof fragment for the phase-two explicit finite-algebra construction.
\section{An explicit polynomial for a guarded finite algebra}
\label{sec:explicit-finite-algebra}

Let $R=\mathbb Z[x_1,\ldots,x_b]$. Consider a triangular monic presentation
\[
 A=R[z_1,\ldots,z_m]/(f_1,\ldots,f_m),\qquad
 f_i\in R[z_1,\ldots,z_i],
\]
where $f_i$ is monic of degree $e_i\ge1$ in $z_i$. Its standard
monomials $\prod_i z_i^{a_i}$, $0\le a_i<e_i$, form an $R$-basis.
In particular $A$ is free of rank $n=\prod_i e_i$. For $a\in A$,
write $M_a$ for its multiplication matrix in this ordered monomial basis.
The matrices $M_a$ commute, and their entries are integer polynomials in
the base coordinates. They are completely specified by repeated reduction
using the monic equations; they need not be materialized to specify a
determinant involving them.

Fix elements $G,E_1,\ldots,E_r\in A$. We assume that in every geometric
fiber, every point at which $G$ is nonzero has local length one. This
condition holds, for example, if $G$ contains the product of the diagonal
derivatives $\partial f_i/\partial z_i$ as a factor. We seek the rational
base points admitting a rational point of $A$ with
\[
 G\ne0,\qquad E_1=\cdots=E_r=0.
\]

For $0\le\ell\le m n(n+1)/2$, define
\[
 \lambda_\ell=1+\sum_{i=1}^m\ell^i z_i,\qquad
 M=M_{G\lambda_\ell},\qquad
 H(V)=\det(I-VM)=\sum_{k=0}^n a_kV^k.
\]
Here and below $a_k\in R$ depend on the fixed integer $\ell$; always
$a_0=1$. Also put
\[
 R_i(V)=\operatorname{Tr}\bigl(M_{E_i}\operatorname{adj}(I-VM)\bigr).
\]
This is a polynomial in $R[V]$ of degree at most $n-1$.

For each $1\le d\le n$, set
\begin{align*}
 H_d(V)&=\sum_{k=0}^d a_kV^k,\\
 g&=a_d\operatorname{Res}_{V}(H_d,H'_d),\qquad L=a_d^2g,\\
 \Psi(U)&=1+\sum_{j=1}^d a_{d-j}a_d^{2j-1}g^jU^j,\\
 \widetilde R_i(U)&=(LU)^{n-1}R_i\bigl(1/(LU)\bigr),\\
 Q_{\ell,d}(\mathbf x,U)&=
 \Psi(U)^2+\sum_{k=d+1}^n a_k^2+\sum_{i=1}^r\widetilde R_i(U)^2.
\end{align*}
The resultant is the Sylvester determinant with the specified degree
bounds $d$ and $d-1$. The expression for $\widetilde R_i$ denotes
polynomial coefficient reversal, hence defines a polynomial also at
$LU=0$. Every displayed expression has integer coefficients.

\begin{theorem}\label{thm:explicit-guarded-finite-algebra}
The explicitly specified polynomial
\[
 P(\mathbf x,U)=
 \prod_{\ell=0}^{m n(n+1)/2}\ \prod_{d=1}^{n} Q_{\ell,d}(\mathbf x,U)
\]
has the following property. For every $\mathbf a\in\mathbb Q^b$,
there exists $U\in\mathbb Q$ with $P(\mathbf a,U)=0$ if and only if
there is a rational point of the fiber algebra $A_{\mathbf a}$ satisfying
$G\ne0$ and all $E_i=0$. The formula uses only one coordinate in addition
to the $b$ base coordinates; its formal algebra generators are eliminated
by the determinants and are not quantified coordinates.
\end{theorem}
\begin{proof}
Fix a rational base point. First consider one factor $Q_{\ell,d}$.
At $a_dg=0$ or $U=0$, its reciprocal residual $\Psi$ equals one.
Otherwise, with $V=1/(LU)$,
\[
 \Psi(U)=(LU)^d\frac{H_d(V)}{a_d}.
\]
Thus a rational zero of $Q_{\ell,d}$ implies $a_dgU\ne0$, all
coefficients $a_k$ with $k>d$ vanish, and $H(V)=H_d(V)=0$.
The resultant condition makes this a simple root of $H$. Since
$H(0)=1$, the root $V$ is nonzero.

The rational matrix $I-VM$ has a one-dimensional kernel. Every
multiplication matrix $M_{z_i}$ preserves it and therefore acts on
it by a rational scalar $\zeta_i$. Applying the algebra relations
to this common eigenline shows that $(\zeta_1,\ldots,\zeta_m)$ is
a rational point of the fiber. Its $G\lambda_\ell$-value is $1/V$,
so its $G$-value is nonzero.

To verify the closed conditions, write
$J=\operatorname{Tr}(\operatorname{adj}(I-VM))$. At a root of $H$,
the adjugate identity gives
$M\operatorname{adj}(I-VM)=V^{-1}\operatorname{adj}(I-VM)$.
Differentiating the determinant then yields
\[
 J=-V H'(V)\ne0.
\]
The columns of the adjugate lie in the common eigenline, so
\[
 R_i(V)=E_i(\boldsymbol\zeta)J.
\]
Since $LU\ne0$, the equation $\widetilde R_i(U)=0$ is equivalent
to $R_i(V)=0$, and hence $E_i(\boldsymbol\zeta)=0$. This proves
the reverse implication without assuming that the rest of the fiber
algebra is reduced.

Conversely, suppose a suitable rational point exists. Over an
algebraic closure, multiplication by $G\lambda_\ell$ has the
value $G(p)\lambda_\ell(p)$ at every geometric point $p$, with
multiplicity equal to its local length. If $G(p)=0$, its factor
in $\det(I-VM)$ is one, even if the local algebra is nonreduced.
If $G(p)\ne0$, its local length is one by hypothesis.

At any such point $p$, the polynomial
$G(p)(1+\sum_{i=1}^m\ell^i z_i(p))$ is nonzero as a polynomial
in $\ell$, because its constant coefficient is $G(p)\ne0$.
For distinct such points $p,q$, their difference is also nonzero:
if all coefficients vanished, first $G(p)=G(q)\ne0$, and then
all their coordinates would agree. Each of these polynomials has
degree at most $m$. With at most $n$ geometric points, there are
at most $m n(n+1)/2$ bad values of $\ell$ in total, counting zero
values and collisions. Thus one of the specified integers makes
all nonzero guarded eigenvalues distinct and nonzero.

For that $\ell$, the polynomial $H$ is squarefree of some degree
$1\le d\le n$. The given rational point provides a nonzero
rational root $V$ of $H$. Its closed residuals vanish, so the trace
identities give $R_i(V)=0$. The leading coefficient and resultant
are nonzero. Taking $U=1/(LV)$ therefore gives a zero of the
corresponding factor $Q_{\ell,d}$, and hence of $P$.
\end{proof}

The degree cases are part of the literal finite formula. They include
every exceptional drop in degree of $H$. The discriminant condition
excludes bad separating choices, and the finite bound on $\ell$ proves
that this exclusion loses no valid rational point. In particular the
construction does not leave a generic primitive-element argument or an
unperformed exceptional-stratum decomposition implicit.
